\clearpage\begingroup\SourceLeaf{350b}\ArithmeticHead{19}{R}{DIE GRUNDLAGEN DER GEOMETRIE.}
\renewcommand{\thefootnote}{\arabic{footnote})}\setcounter{footnote}{0}
\begin{center}\textbf{4.}\par\smallskip
\textls[100]{Fortschritte in den Grundlagen der Geometrie (1795—1799).}
\end{center}
\noindent\hspace*{\parindent}Als \textsc{Gauss} im Oktober 1795 seine Studien in Göttingen begann, hatte
er bereits, wie wir bemerkten, die schwache Stelle des \textsc{Euklid}ischen Lehrgebäudes
erkannt und war wenigstens bei dem Inhalt der Vielecke den Folgerungen
nachgegangen, die sich aus der Verwerfung des Parallelenaxioms ergeben.
\par
Untersuchungen über die Grundlagen der Geometrie haben gegen das
Ende des achtzehnten Jahrhunderts die Mathematiker und darüber hinaus weite
Kreise der Gebildeten lebhaft beschäftigt. Von zwei Seiten waren Anregungen
dazu gekommen.
\par
Seit dem Ende des 17. Jahrhunderts hatten, um nur einige wichtige
Namen zu nennen, \textsc{Hume}, \textsc{Leibniz}, \textsc{d’Alembert} die Frage nach dem Wesen
der mathematischen Erkenntnis aufgeworfen, und durch \textsc{Kant}s \textit{Kritik der
reinen Vernunft} (1781, 1787) war diese Frage geradezu in den Mittelpunkt
der philosophischen Erörterungen gestellt worden. Dabei war es besonders
die Parallelentheorie, an der sich Berufene und Unberufene versuchten, denn
dem elften \textsc{Euklid}ischen Axiom fehlte jenes Merkmal der einleuchtenden Gewissheit,
die dem Apriorischen eigen sein sollte\footnote{I. \textsc{Kant}, \textit{Kritik der reinen Vernunft}, 1. Aufl. 1781, S. 25, 2. Aufl. 1777, S. 39: »So werden auch alle geometrischen Grundsätze .\ .\ .\ niemals aus allgemeinen Begriffen .\ .\ ., sondern aus der Anschauung, und zwar \textls[100]{a priori} mit apodiktischer Gewissheit hergeleitet«.}; es ist deutlich zu erkennen,
\par\vfill\hfill{\footnotesize 3\textsuperscript{*}}\par
\clearpage\SourceLeaf{351}\ArithmeticHead{20}{L}{STÄCKEL, GAUSS ALS GEOMETER.}\setcounter{footnote}{0}
\noindent wie mit dem Jahre 1781 die Flut der Veröffentlichungen anschwillt, die sich
auf die Parallelenfrage beziehen\footnote{Vgl. das Literaturverzeichnis bei P. \textsc{Stäckel} und F. \textsc{Engel}, \textit{die Theorie der Parallellinien von Euklid bis auf Gauss}, Leipzig 1895; im Folgenden angeführt mit P. Th.}.
\par
Noch von einer anderen Seite kamen Einwirkungen. Der französische
Umsturz führte zur Gründung neuer Hochschulen in Paris, der École polytechnique
und der École normale, an die die bedeutendsten Mathematiker des
Landes berufen wurden. Dies veranlasste sie, zu den Elementen ihrer Wissenschaft
zurückzukehren. \textsc{Lagrange} verschmähte es nicht, Vorlesungen über
die Elemente der Arithmetik und Algebra zu halten\footnote{J. L. \textsc{Lagrange}, \textit{Leçons élémentaires sur les mathématiques, données à l’École normale en 1795}, Oeuvres t. 7, S. 183.}, und \textsc{Legendre} liess
1794 seine Elemente der Geometrie erscheinen, die einen ungewöhnlichen
Erfolg hatten und 1823 ihre zwölfte Auflage erlebten\footnote{A. M. \textsc{Legendre}, \textit{Eléments de géométrie}, Paris 1794; 12. éd., Paris 1823.}. In der Parallelentheorie
war \textsc{Legendre} bemüht gewesen, die bei \textsc{Euklid} vorhandenen Mängel
zu beseitigen, aber die beständigen Änderungen bei den auf einander folgenden
Auflagen zeigen, dass er auf schwankendem Boden stand; wie seine letzte,
zusammenfassende Veröffentlichung vom Jahre 1833 erkennen lässt\footnote{A. M. \textsc{Legendre}, \textit{Réflexions sur les différentes manières de démontrer la théorie des parallèles}, Mém. de l’Acad., t. 12, année 1828, Paris 1833, S. 367.}, hat er
sich niemals zu dem Gedanken der Unbeweisbarkeit des elften Axioms erheben
können\footnote{In dem Briefe an \textsc{Olbers} vom 30. Juli 1806 (W. VIII, S. 139, 165) bemerkt \textsc{Gauss}, es scheine sein
Schicksal zu sein, in fast allen seinen theoretischen Arbeiten mit \textsc{Legendre} zu konkurrieren, und führt
dafür an: die höhere Arithmetik, die transzendenten Funktionen, welche mit der Rektifikation der Ellipse
zusammenhängen, die ersten Gründe der Geometrie und die Methode der kleinsten Quadrate.
\par
Für die höhere Arithmetik kommt in Betracht \textsc{Legendre}s \textit{Essay sur la théorie des nombres}, Paris
1798, dessen Verhältnis zu den \textit{Disquisitiones arithmeticae} \textsc{Tschebyscheff} in seiner \textit{Theorie der Kongruenzen}
(deutsch von \textsc{Schapira}, Berlin 1889) gut gekennzeichnet hat; im Besonderen ist noch das Reziprozitätsgesetz
der quadratischen Reste zu nennen; vgl. \textsc{Bachmann}, W. X\textsubscript{2}, Abh. I, S. 14. Die elliptischen Integrale hat
\textsc{Legendre} in dem grundlegenden \textit{Mémoire sur les transcendantes elliptiques}, Paris 1794 behandelt und ihnen
dann zwei umfangreiche Werke gewidmet: \textit{Exercices de calcul intégral}, 3 Bände, Paris 1811—1816; \textit{Traité
des fonctions elliptiques}, 3 Bände, Paris 1825—1832. Die Methode der kleinsten Quadrate entwickelt \textsc{Legendre}
in den \textit{Nouvelles méthodes pour la détermination des orbites des comètes}, Paris 1805, während die
\textit{Theoria motus corporum coelestium} von \textsc{Gauss} erst 1809 erschienen ist (vgl. auch W. VIII, S. 136—141
und X\textsubscript{1}, S. 373 und 380).
\par
Hinzuzufügen wäre noch, dass \textsc{Gauss} und \textsc{Legendre} sich mit der Theorie und Praxis der Geodäsie
beschäftigt haben und dass \textsc{Legendre}s Satz über die Zurückführung eines kleinen sphärischen Dreiecks
auf ein ebenes Dreieck mit ebenso langen Seiten auf die Untersuchungen von \textsc{Gauss} zur allgemeinen Lehre
von den krummen Flächen anregend gewirkt hat. Auch bei der Anziehung der homogenen Ellipsoide sind
beide zusammengetroffen; für \textsc{Legendre} sind hier zu nennen die Abhandlungen in den Mémoires des savants
étrangers, t. 10, Paris 1785 und in den Mémoires de l’Institut, année 1810, 2. partie, Paris 1814. Endlich
sind noch die Arbeiten über das von \textsc{Gauss} mit \(\Pi\), von \textsc{Legendre} mit \(\Gamma\) bezeichnete \textsc{Euler}sche Integral
zu erwähnen (\textit{Exercices}, t. I, S. 222—307).
\par
Die Vergleichung der Leistungen zeigt, dass \textsc{Legendre} mit scharfem Blick die Stellen erkannt hatte,
an denen die mathematische Forschung mit Erfolg einsetzen konnte. Seinem unermüdlichen Fleiss und
analytischen Geschick ist eine Reihe schöner Erfolge zu Teil geworden, jedoch blieb er überall auf einer
Stufe stehen, die zu überschreiten erst dem Genie von \textsc{Gauss} vergönnt war. Mit besonderer Deutlichkeit
tritt dies bei den Grundlagen der Geometrie hervor.}.
\clearpage\SourceLeaf{352}\ArithmeticHead{21}{R}{DIE GRUNDLAGEN DER GEOMETRIE.}\setcounter{footnote}{0}
\noindent\hspace*{\parindent}Die Universität Göttingen nahm lebhaften Anteil an der Bewegung, deren
Hervortreten soeben geschildert wurde. Professor der Mathematik war damals
\textsc{Kaestner} (1719—1800). Er hat die Literatur über die Parallelentheorie eifrig
gesammelt und eine noch heute wertvolle Dissertation, \textsc{Klügel}s \textit{Recensio conatuum
praecipuorum theoriam parallelarum demonstrandi} vom Jahre 1763, veranlasst.
In dem Nachwort meint \textsc{Kaestner}, ein Beweis des Parallelenaxioms
sei nur zu erhoffen durch eine genauere Ausbildung der Geometrie der Lage,
die mit \textsc{Leibniz} untergegangen sei. Gegenwärtig bleibe nur übrig, offen die
Forderung \textsc{Euklid}s als solche auszusprechen; niemand, der bei gesunden Sinnen
sei, werde sie bestreiten wollen. In seinen späteren Vorlesungen hat \textsc{Kaestner}
»an der Möglichkeit der Lösung verzweifelnd mit unbegreiflicher Resignation,
anstatt nach der wahren Demonstration zu forschen, ein blindes Annehmen
angeraten« (P. Th. S. 139—141). Ähnlich wie \textsc{Kaestner} dachte auch
sein Kollege an der Nachbar-Universität Helmstedt, \textsc{Joh. Friedr. Pfaff} (1765—1825),
der meinte, alles was sich tun liesse sei, das Parallelenaxiom durch ein
einfacheres zu ersetzen, es zu simplifizieren (P. Th. S. 215).
\par
Als \textsc{Gauss} nach Göttingen kam, habilitierte sich gerade für Mathematik
J. \textsc{Wildt} (1770—1844) mit einer Probeschrift über die Parallelentheorie\footnote{J. \textsc{Wildt}, \textit{Theses quae de lineis parallelis respondent}, Göttingen 1795. \textsc{Wildt} hat in den Göttinger Gelehrten Anzeigen, Jahrgang 1800, S. 1769—1772 drei »auf reiner Anschauung beruhende Beweise« des elften Axioms veröffentlicht.}.
Ein Liebhaber auf diesem Gebiete war auch der ausserordentliche Professor
der Astronomie \textsc{Carl Felix Seyffer} (1762—1822). Im Jahre 1801 hat er
zwei Besprechungen von Versuchen, das Parallelenaxiom zu beweisen, in den
Göttinger Gelehrten Anzeigen veröffentlicht; sie zeigen, dass er die Schriften
mit Verständnis und Urteil gelesen hatte; ja \textsc{Seyffer} war zu der Einsicht
\clearpage\SourceLeaf{353}\ArithmeticHead{22}{L}{STÄCKEL, GAUSS ALS GEOMETER.}\setcounter{footnote}{0}
\noindent gekommen, dass »es mehr als zweifelhaft scheine, ob es überhaupt möglich
sei, das elfte Axiom zu beweisen, ohne ein neues Axiom zu Hilfe zu nehmen«
(P. Th. S. 214).
\par
Während \textsc{Gauss} zu \textsc{Kaestner} und \textsc{Wildt} in kein näheres Verhältnis getreten
ist, hat er mit \textsc{Seyffer} verkehrt und ist mit ihm bis zu dessen Tode
im Briefwechsel geblieben. Ihre Unterhaltungen haben auch die Parallelentheorie
betroffen, denn am 26. Juni 1801 schreibt \textsc{Seyffer} an \textsc{Gauss}: »Vielleicht
ist es Ihnen nicht uninteressant, dass die Rezension in der hiesigen
Zeitung [den Gött. Gel. Anzeigen] über die Theorie der Parallelen von \textsc{Schwab}
von mir war. Ich wünschte, dass Sie mir Ihre lehrreichen Ideen hierüber
gelegentlich sagten« (Brief im \textsc{Gauss}-Archiv).
\par
Im Hause \textsc{Seyffer}s hat \textsc{Gauss} seinen besten Jugendfreund, den Ungarn
\textsc{Wolfgang Bolyai} kennen gelernt. »Als \textsc{Wolfgang} nach Göttingen kam«,
erzählt sein Sohn \textsc{Johann}, »traf er mit \textsc{Gauss} zufällig bei dem Professor
[\textsc{Seyffer}] zusammen und äusserte sich da freimütig und entschieden über die
Leichtfertigkeit der Behandlung der Mathematik; kurz darauf begegnete er
\textsc{Gauss} am Walle beim Spazierengehen; sie näherten sich einander. Mein
Vater sprach unter anderem von seinen Gedanken behufs Erklärung der geraden
Linie und der etwaigen Wege zum Beweise des elften Axioms, und
der damals schon zum Koloss in den höheren Regionen der Wissenschaft,
besonders der Zahlenlehre, emporgewachsene \textsc{Gauss} brach ergötzt, überrascht
in die lakonischen Worte aus: Sie sind ein Genie; Sie sind mein Freund!,
worauf sogleich das Band der Brüderschaft erfolgte«\footnote{\textsc{Wolfgang} und \textsc{Johann Bolyai}, \textit{Geometrische Untersuchungen}, herausgegeben von P. \textsc{Stäckel}, Leipzig 1913, 1. Teil: Leben und Schriften der beiden \textsc{Bolyai}, S. 8; im Folgenden angeführt mit \textsc{Bol.}}. Über den Verkehr
zwischen den beiden Freunden berichtet \textsc{Wolfgang}: »Er war sehr bescheiden
und zeigte wenig; nicht drei Tage, wie mit Plato, jahrelang konnte man mit
ihm zusammen sein, ohne seine Grösse zu erkennen. Schade, dass ich dieses
titellose, schweigsame Buch nicht aufzumachen und zu lesen verstand. Ich
wusste nicht, wie viel er weiss, und er hielt, nachdem er meine Art sah, viel
von mir, ohne zu wissen, wie wenig ich bin. Uns verband die wahre (nicht
oberflächliche) Leidenschaft für die Mathematik und unsere sittliche Übereinstimmung,
so dass wir oft, mit einander wandernd, mit den eigenen Gedanken
beschäftigt stundenlang wortlos waren« (\textsc{Bol.} S. 9).
\clearpage\SourceLeaf{354}\ArithmeticHead{23}{R}{DIE GRUNDLAGEN DER GEOMETRIE.}\setcounter{footnote}{0}
\noindent\hspace*{\parindent}Was \textsc{Bolyai} und \textsc{Gauss} über das Parallelenaxiom mit einander verhandelt
haben, wissen wir nicht. Wohl aber wissen wir, dass \textsc{Wolfgang}, nachdem
sein Freund im Herbst 1798 nach Braunschweig zurückgekehrt war, sich angestrengt
bemüht hat, das Axiom zu beweisen und dass er im Mai 1799 sein
Ziel erreicht zu haben glaubte. \textsc{Bolyai} ist nämlich, ehe er Deutschland verliess,
noch einmal mit \textsc{Gauss} zusammengetroffen. Am 24. Mai 1799 haben
die beiden zu Klausthal im Harz von einander Abschied genommen, und bei
dieser Zusammenkunft hat \textsc{Wolfgang} von seiner »Göttinger Parallelentheorie«
erzählt. Hierauf bezieht sich eine Stelle des Briefes von \textsc{Gauss} an \textsc{Wolfgang}
vom 16. Dezember 1799: »Es tut mir sehr leid, dass ich unsere ehemalige
grössere Nähe nicht benutzt habe, um mehr von Deinen Arbeiten über die
ersten Gründe der Geometrie zu erfahren; ich würde mir gewiss dadurch
manche vergebliche Mühe erspart haben und ruhiger geworden sein, als jemand
wie ich es sein kann, so lange bei einem solchen Gegenstande noch so viel
zu desiderieren ist. Ich selbst bin in meinen Arbeiten darüber weit vorgerückt
(wiewohl mir meine andern ganz heterogenen Geschäfte wenig Zeit dazu
lassen); allein der Weg, den ich eingeschlagen habe, führt nicht so wohl zu
dem Ziele, das man wünscht und welches Du erreicht zu haben versicherst,
als vielmehr dahin, die Wahrheit der Geometrie zweifelhaft zu machen« (W.
VIII, S. 159).
\par
Die Ergebnisse, zu denen \textsc{Gauss}, wie das \textit{Tagebuch} (T. Nr. 99) zeigt, im
September 1799 gelangt war, hat er in dem Briefe nur angedeutet. Er fährt
fort: »Zwar bin ich auf manches gekommen, was den meisten schon für einen
Beweis gelten würde, aber was in meinen Augen so gut wie \textls[100]{Nichts} beweist,
z. B. wenn man beweisen könnte, dass ein geradliniges Dreieck möglich sei,
dessen Inhalt grösser wäre als jede gegebene Fläche, so bin ich im Stande
die ganze Geometrie völlig strenge zu beweisen. Die meisten würden nun
wohl jenes als ein Axiom gelten lassen; ich nicht; es wäre ja wohl möglich,
dass, so entfernt man auch die drei Eckpunkte des Dreiecks im Raume von
einander annähme, doch der Inhalt immer unter (infra) einer gegebenen
Grenze wäre. Dergleichen Sätze habe ich mehrere, aber in keinem finde ich
etwas Befriedigendes« (W. VIII, S. 159).
\par
Aufzeichnungen über die Untersuchungen, von denen \textsc{Gauss} spricht, sind
uns nicht erhalten. Es ist jedoch sehr wahrscheinlich, dass der Brief an
\clearpage\SourceLeaf{355}\ArithmeticHead{24}{L}{STÄCKEL, GAUSS ALS GEOMETER.}\setcounter{footnote}{0}
\noindent\textsc{Bolyai} vom 6. März 1832 (W. VIII, S. 220) einen Teil dieser Untersuchungen
wiedergibt. In diesem wiederholt angeführten Briefe sagt \textsc{Gauss}, dass er schon vor
30 bis 35 Jahren Meditationen über die Grundlagen der Geometrie angestellt habe
und dass zu seiner Überraschung die Ergebnisse der \textit{Scientia spatii} \textsc{Johann Bolyai}s
fast durchgehends damit übereinstimmten. In manchem Teile habe er etwas
andere Wege eingeschlagen und als ein Specimen füge er in den Hauptzügen
einen rein geometrischen Beweis des Lehrsatzes bei, dass in der antieuklidischen
Geometrie die Differenz der Winkelsumme eines Dreiecks von \(180^0\) dem
Flächeninhalte proportional sei. Der Beweis beginnt mit dem Satze, dass das
asymptotische Dreieck, bei dem die drei Ecken im Unendlichen liegen, eine
bestimmte endliche Area habe. Eine Herleitung wird nicht angegeben\footnote{Man kann durch eine einfache, nur die allerersten Eigenschaften asymptotischer Geraden benutzende Zeichnung ein solches Dreieck in ein inhaltgleiches, ganz im Endlichen liegendes Viereck verwandeln. Vgl. H. \textsc{Liebmann}, \textit{Zur nichteuklidischen Geometrie}, Leipziger Berichte, Bd. 58, 1906, S. 560; \textit{Nichteuklidische Geometrie}, 2. Aufl., Leipzig 1912, S. 53.}. Für
den Inhalt eines Dreiecks, bei dem eine Ecke im Endlichen liegt, während
die Gegenseite zu den beiden anderen Seiten asymptotisch ist, ergibt sich
dann eine Funktionalgleichung, die im Gebiete der stetigen Funktionen leicht
gelöst werden kann. Nun entsteht ein ganz im Endlichen liegendes Dreieck
aus einem asymptotischen Dreieck durch Wegnahme von solchen Dreiecken,
bei denen eine Ecke im Endlichen liegt, und so folgt schliesslich die zu beweisende
Behauptung.
\par
Wie immer auch \textsc{Gauss} im Jahre 1799 vorgegangen sein mag, so zeigt
sein Brief von 16. Dezember 1799 auf jeden Fall, dass er sich damals auf
dem Wege befand, den vor ihm \textsc{Saccheri} (1733) und \textsc{Lambert} (1766) eingeschlagen
hatten, nämlich planmässig die Folgerungen zu entwickeln, die sich
aus der Annahme ergeben, das \textsc{Euklid}ische Parallelenaxiom sei nicht erfüllt.
Da \textsc{Gauss} hierbei auf keinen Widerspruch kam, wurde ihm die Wahrheit der
\textsc{Euklid}ischen Geometrie zweifelhaft. Den Gedanken, dass die nichteuklidische
Geometrie »wahr« sein könne, hatte übrigens schon \textsc{Lambert} offen ausgesprochen
(P. Th. S. 200).
\par
Hierbei erheben sich die Fragen, ob \textsc{Gauss} jene Arbeiten gekannt und
wann er sie möglicher Weise kennen gelernt hat. Gewiss sind sie ihm in
der Göttinger Universitätsbibliothek, die er als Student fleissig benutzt hat,
\clearpage\SourceLeaf{356}\ArithmeticHead{25}{R}{DIE GRUNDLAGEN DER GEOMETRIE.}\setcounter{footnote}{0}
\noindent zugänglich gewesen. Allein man muss bedenken, dass diese Schriften \textsc{Wolfgang
Bolyai} unbekannt geblieben sind; dies geht mit voller Sicherheit aus
den Äusserungen seines Sohnes hervor (\textsc{Bol.} S. 221—223).
\par
Dass andererseits später in den Kreisen der Schüler von \textsc{Gauss} von \textsc{Lambert}s
Theorie der Parallellinien gesprochen wurde, zeigen Briefe von \textsc{Bessel}
an \textsc{Encke} vom 9. Juli 1821 und von \textsc{Encke} an \textsc{Bessel} vom 13. Oktober 1821
(Abschriften im \textsc{Gauss}-Archiv). Auch wird \textsc{Lambert} in dem Briefe \textsc{Bessel}s
an \textsc{Gauss} vom 10. Februar 1829 erwähnt (W. VIII, S. 201). Endlich besass
\textsc{Gauss} die \textit{Mathematischen Abhandlungen} von J. W. H. \textsc{Lehmann}, Zerbst 1829,
in denen \textsc{Saccheri} und \textsc{Lambert} angeführt werden; Randbemerkungen und
Spuren des Gebrauches lassen darauf schliessen, dass \textsc{Gauss} darin gelesen und
die auf die Parallelentheorie bezüglichen Stellen beachtet hat\footnote{Vgl. den Aufsatz von P. \textsc{Stäckel}: \textit{F. A. Taurinus}, Abhandlungen zur Geschichte der Mathematik, Heft 9, Leipzig 1899, S. 427.}.
\par
Entscheidend für die Beurteilung der Leistung von \textsc{Gauss} ist der Umstand,
dass weder \textsc{Saccheri} noch \textsc{Lambert} bis zur nichteuklidischen ebenen Trigonometrie
vorgedrungen sind, wenn ihr auch \textsc{Lambert} durch den Gedanken, »die
dritte Hypothese komme bei einer imaginären Kugelfläche vor« (P. Th. S. 203)
nahe gekommen war; denn erst die Trigonometrie sichert für die Ebene die
Widerspruchslosigkeit der absoluten Geometrie und führt damit zu der Überzeugung,
dass alle Versuche, das Parallelenaxiom durch Konstruktionen in
der Ebene zu beweisen, vergeblich sein müssen.
\clearpage\endgroup
